\chapter{Theoretical Background} \label{ch:theoretical_background}

\section{Cosmological Background} \label{sec:cosmological_background}
The foundation of modern cosmology is based on the \textit{Cosmological Principle}, which states that the Universe 
is homogeneous and isotropic on large scales. Homogeneity implies that the large-scale properties of the Universe 
are statistically uniform across different locations, whereas isotropy means that no direction is preferred. Although 
this principle is introduced as a fundamental assumption in modern cosmology, it also represents a hypothesis that can be 
tested through observations. Current astronomical evidence strongly supports the validity of the cosmological principle on the 
largest observable scales (e.g.~\cite{Scrimgeour_2012, Laurent_2016})

Based on the cosmological principle, the Universe can be described by the \textit{Friedmann-Lemaître-Robertson-Walker} (FLRW) metric

\begin{equation*}
    ds^2 = -c^2 dt^2 +a^2(t) \left[ \frac{dr^2}{1 - k r^2} + r^2 (d\theta^2 + \sin^2\theta d\phi^2) \right]
\end{equation*}

with the spatial comoving coordinates $(r, \theta, \phi)$, the speed of light $c$, and the cosmic time $t$~\cite{CervantesCota_2011}. The parameter $k$ describes the curvature of the Universe, 
which can take values of $-1$, $0$, or $+1$ for open, flat, and closed geometries, respectively. 
The scale factor $a(t)$ describes the relative expansion of the Universe and is conventionally normalized to its present value,

\begin{equation*}
    a(t_0)=1.
\end{equation*}

For earlier cosmic times, the scale factor satisfies $a(t)<1$, whereas in an expanding Universe it increases with time and can exceed unity in the future. 
A static Universe would correspond to $\dot{a}(t)=0$, resulting in a vanishing Hubble parameter~\cite{Huterer_2018}.

The expansion rate of the Universe is quantified by the Hubble parameter,

\begin{equation*}
    H(t)=\frac{\dot{a}(t)}{a(t)},
\end{equation*}

which describes the relative change of the scale factor with cosmic time. 
At the present time

\begin{equation*}
    H_0 = H(t_0)
\end{equation*}

is referred to as the Hubble constant~\cite{CervantesCota_2011}.

The expansion of space introduces a distinction between physical distances and 
comoving distances. While the physical distance between two objects changes 
with time due to the expansion of the Universe, the comoving distance remains 
constant for objects that follow the general expansion of space. The physical or so called
proper distance is related to the comoving distance $\chi$ by

\begin{equation*}
    D(t)=a(t) \chi.
\end{equation*}

Thus, comoving coordinates remove 
the effect of the cosmic expansion and provide a convenient description of the 
large-scale structure of the Universe~\cite{Ishak_2019}.

The expansion of the Universe can be observed through the cosmological 
redshift. Due to the expansion of space, the wavelength of photons emitted by 
distant objects is stretched during their propagation. 
The redshift is defined as

\begin{equation*}
    z=\frac{\lambda-\lambda_0}{\lambda_0},
\end{equation*}

where $\lambda_0$ is the emitted wavelength and $\lambda$ is the observed 
wavelength. The redshift is directly related to the scale factor through

\begin{equation*}
    1+z = \frac{a(t_0)}{a(t_e)}
\end{equation*}

with $t_e$ denoting the time of emission.
Consequently, the observed redshift provides information about the expansion 
history of the Universe and can be used to determine cosmological distances. 
The line-of-sight comoving distance to an object at redshift $z$ is given by

\begin{equation*}
    \chi(z)= \chi_H \int_0^z \frac{\mathrm{d}z'}{E(a(z'))},
\end{equation*}

where $\chi_H = c/H_0$ is the Hubble distance and $E(a)=H(a)/H_0$ is the dimensionless Hubble parameter~\cite{Hogg_2000}. 


\section{Friedmann Equations and Cosmic Expansion} \label{sec:friedmann_equations_and_cosmic_expansion}
Following the cosmological principle, the matter and energy content of the Universe can be described by 
a perfect-fluid energy-momentum tensor. In combination with the FLRW metric, Einstein's field equations 
reduce to the \textit{Friedmann equations}

\begin{equation*}
    H^2 = \frac{\dot{a}^2}{a^2} = \frac{8 \pi G}{3} \rho - \frac{k c^2}{a^2} + \frac{\symup{\Lambda}}{3}
\end{equation*}

\begin{equation*}
    \frac{\ddot{a}}{a} = -\frac{4 \pi G}{3} \left( \rho + \frac{3 p}{c^2} \right) + \frac{\symup{\Lambda}}{3}
\end{equation*}

which determine the evolution of the scale factor $a(t)$. Here $\rho$ and $p$ denote the energy density and pressure of the cosmic fluid, 
$\symup{\Lambda}$ is the cosmological constant, and $G$ is the gravitational constant~\cite{Huterer_2018}.

For a spatially flat Universe ($k=0$), the critical density is defined as the 
energy density required to obtain the observed expansion rate $H$,

\begin{equation*}
    \rho_c=\frac{3H^2}{8\pi G}.
\end{equation*}

The ratio between the density of a cosmological component and the critical 
density defines the corresponding density parameter,

\begin{equation*}
    \symup{\Omega_i}=\frac{\rho_i}{\rho_c}.
\end{equation*}

A flat Universe therefore satisfies

\begin{equation*}
    \symup{\Omega}_{0}= \symup{\Omega}_m + \symup{\Omega}_r + \symup{\Omega}_{\symup{\Lambda}} + \symup{\Omega}_k = 1,
\end{equation*}

where the total density includes matter, radiation, dark energy, and possible 
curvature contributions.

Consequently, the Friedmann equation can be expressed in terms of the density parameters~\cite{BartelmannSchneider_2001}:

\begin{equation*}
    H^2 = H_0^2 \left[ \symup{\Omega}_m a^{-3} + \symup{\Omega}_r a^{-4} + \symup{\Omega}_{\symup{\Lambda}} + (1-\symup{\Omega}_{0}) a^{-2} \right].
\end{equation*}

HIER NOCH KLAEREN WIESO GLEICH 1???
HIER SO SACHEN KLAEREN WIE WAS GENAU IST DICHTE UND DRUCK,EINFACH VERSTEHEN UND IN NOTIZEN MARKIEREN

\section{\texorpdfstring{The $\symup{\Lambda}$CDM Model}{The Lambda CDM Model}} \label{sec:lcdm_model}
HIER VOR ALLEM SCHAUEN WAS IM CHAT FENSTER ALLES STEHT (WAHRESCHEINLICH ZU AUSFUEHRELICH DORT)SIEHE BACHELORARBEIT BERTMANN

\begin{itemize}
    \item Matter
    \item Radiation
    \item Dark energy
    \item $H(z)$
    \item Cosmological parameters
\end{itemize}


\section{Fast Radio Bursts} \label{sec:fast_radio_bursts}
\begin{itemize}
 \item Discovery and Properties
 \item Physical Origin
 \item Cosmological Applications
 \item magnetar and neutron star as progenitors , in which areas they occur
\end{itemize}

\section{Galaxy Surveys and Catalogs} \label{sec:galaxy_surveys_and_catalogs}
\begin{itemize}
 \item Galaxy Survey DR5 stage 3 --> stage 54 KiDS, how measured
 \item COSMOS2020 Catalog ?
 \item Comparison of the Data Sets
\end{itemize}

\section{Correlation Functions} \label{sec:correlation_functions}
\begin{itemize}
 \item Two-Point Correlation Function
 \item Auto-Correlation
 \item Cross-Correlation
 \item how do you get to Cell
\end{itemize}